% Part 1, Problem 16: Density of rationals - Existence Proof
% Hard

\begin{problem}[Hard]
Prove that for any real numbers $x$ and $y$ with $x < y$, there exists a rational number $q$ such that
\[
x < q < y.
\]
\end{problem}

\begin{solution}
\textbf{Existence Proof}

Since $y-x>0$, we can choose a positive integer $n$ such that
\[
n(y-x)>1.
\]
Equivalently,
\[
ny-nx>1.
\]

Now let $m$ be the smallest integer greater than $nx$. Then
\[
m-1 \le nx < m.
\]
So in particular,
\[
nx < m \le nx+1.
\]

But $ny-nx>1$ gives
\[
nx+1<ny.
\]
Hence
\[
nx<m\le nx+1<ny,
\]
so
\[
nx<m<ny.
\]

Dividing through by $n>0$,
\[
x<\frac{m}{n}<y.
\]

Therefore, with $q=\frac{m}{n}$, we have found a rational number satisfying $x<q<y$.

\hfill $\square$
\end{solution}

\begin{takeaways}
\begin{itemize}
\item \textbf{Magnification Idea:} Multiply the interval by a large integer so its length becomes greater than $1$
\item \textbf{Key Step:} Once the gap exceeds $1$, an integer can be placed between the endpoints
\item \textbf{Density of $\mathbb{Q}$:} Rational numbers occur between every two distinct real numbers
\item \textbf{Big Picture:} This is an existence proof rather than a constructive formula for a specific rational
\item For any real numbers $x$ and $y$ with $x < y$, there exists infinitely many rational numbers $q$ such that $x < q < y$.

\end{itemize}
\end{takeaways}
